% TOP_PROBLEM: {"id": 3100021, "title": "\"Seymour's Second Neighborhood Conjecture\" Any oriented graph has a vertex whose outdegree is at most its second outdegr", "category_id": 2, "category": {"id": 2, "name": "combinatorics", "display_name": "Combinatorics", "description": "Counting problems, graph theory, discrete structures.", "slug": "combinatorics", "order_index": 2, "created_at": "2026-07-31T15:26:25.671Z"}, "status": "partially_solved", "problem_number": "AMR-030-0021", "set_id": 13}
% FIRST_POSTED: 2026-10-02
% ENTRY_KIND: statement_only
% UnsolvedMath ID 3100021; source catalogue code AMR-030-0021
% !TeX program = lualatex
% Definition and sources only; this file is not an LLM solution attempt.
\documentclass[11pt,a4paper]{article}
\usepackage[margin=25mm]{geometry}
\usepackage{fontspec}
\setmainfont{lmroman10-regular.otf}[BoldFont=lmroman10-bold.otf,
 ItalicFont=lmroman10-italic.otf,BoldItalicFont=lmroman10-bolditalic.otf]
\usepackage{amsmath,amssymb,mathtools}
\usepackage{unicode-math}
\setmathfont{latinmodern-math.otf}
\AtBeginDocument{\let\setminus\smallsetminus}
\usepackage{xurl,hyperref}
\hypersetup{colorlinks=true,urlcolor=blue}
\setcounter{secnumdepth}{0}
\setlength{\parindent}{0pt}
\setlength{\parskip}{5pt}
\setlength{\emergencystretch}{3em}
\begin{document}
\section*{"Seymour's Second Neighborhood Conjecture" Any oriented graph has a vertex whose outdegree is at most its second outdegr}\label{3100021}
{\small UnsolvedMath ID 3100021 · AMR-030-0021 · Combinatorics · AMR Open Problem Lists}
\subsection{Definitions and mathematical statement}
An oriented graph is a finite directed graph with no loops, no
parallel arcs, and no pair of opposite arcs. For a vertex $v$, put
\[
 N_1^+(v)=\{u:v\to u\},
\]
and let
\[
 N_2^+(v)=\{w\notin N_1^+(v)\cup\{v\}:
                     v\to u\to w\text{ for some }u\}.
\]
The latter is the set of vertices at directed distance exactly
two from $v$. Multiple paths to the same vertex count only once.

Seymour's second neighborhood conjecture asks whether every
nonempty oriented graph has a vertex $v$ satisfying
\[
                         |N_2^+(v)|\ge |N_1^+(v)|.
\]
A sink meets this inequality because both sets are empty, so a
potential counterexample necessarily has positive outdegree at
every vertex. No assumption of strong connectivity is made.

Excluding first neighbors from the second neighborhood matters:
having a two-step walk to an existing first neighbor does not
create a new vertex at distance two. Nor does the question compare
the number of walks of lengths one and two. The orientation
hypothesis rules out directed digons, for which the assertion in
this form can fail. A proof must find at least one suitable vertex;
the inequality is not asserted at every vertex of the graph.
\subsection{Short English statement}
Must a nonempty oriented graph contain a vertex with at least as many distinct second outneighbors as first outneighbors?
\subsection{Sources}
\begin{itemize}
\item[S1] ULAM.AI and the UnsolvedMath Contributors, supplied problems.json, record 3100021 (AMR-030-0021), AMR Open Problem Lists. Catalogue identity and source transcription; CC BY 4.0.
\url{https://huggingface.co/datasets/ulamai/UnsolvedMath}.
\item[S2] Cooper - Combinatorial Problems I Like (2020 snapshot), Graphs, Hypergraphs, Set Systems, Designs, source-order bullet 21. Original source indexed by AMR-030-0021.
\url{https://people.math.sc.edu/cooper/combprob.html}.
\end{itemize}
\end{document}
